% !TeX program = pdflatex
\documentclass[11pt,a4paper]{article}
\usepackage{../../recursos/latex/estilo-notas}

\begin{document}
\cabecalhoaula{06}{Pré-processamento e \emph{Pipelines}}{ISLP labs cap.~5--6; AME §2.1.2}

%==============================================================================
\section*{Objetivos da aula}
%==============================================================================
Ao final desta aula o aluno deve ser capaz de:
\begin{itemize}[leftmargin=1.4cm]
  \item explicar por que o pré-processamento faz parte do \emph{modelo}, e não dos
        dados;
  \item definir a \textbf{padronização} (\texttt{StandardScaler}) e listar suas
        vantagens e desvantagens;
  \item identificar e evitar o \textbf{vazamento de dados} (\emph{data leakage});
  \item construir um \textbf{\emph{pipeline}} que encadeia pré-processamento e
        estimador como um único objeto;
  \item combinar \emph{pipeline} com validação cruzada e busca de hiperparâmetros
        sem vazamento;
  \item conhecer outros pré-processamentos comuns (imputação, codificação de
        categóricas) e o \texttt{ColumnTransformer}.
\end{itemize}

\begin{emsala}
Esta é a aula que costura as anteriores em \emph{boas práticas de implementação}.
Vários métodos já exigiram padronização: Lasso/Ridge (Aula~02) e KNN (Aula~04); e a
Aula~03 já avisou sobre vazamento. Aqui formalizamos \emph{como} fazer isso certo.
Pergunta de abertura: \emph{``onde, exatamente, eu devo calcular a média e o
desvio para padronizar: antes ou depois de separar treino e teste?''}
\end{emsala}

%==============================================================================
\section{Motivação: o pré-processamento é parte do modelo}
%==============================================================================
Na prática, raramente alimentamos o estimador com os dados brutos. Um fluxo típico
encadeia várias transformações: padronizar, imputar faltantes, codificar variáveis
categóricas, reduzir dimensão (PCA), e só então ajustar o modelo. A tentação é
aplicar essas transformações ``de uma vez'' no conjunto inteiro --- e é exatamente aí
que mora o erro.

\begin{ideia}
O insight central da aula: \textbf{toda transformação que ``aprende'' algo dos
dados} (uma média, um desvio, as categorias presentes, os componentes principais)
é \emph{parte do procedimento de estimação}. Logo, deve ser \emph{ajustada apenas
no treino} e \emph{aplicada} ao teste/validação --- nunca o contrário.
\end{ideia}

A ferramenta do \texttt{scikit-learn} para isso é o \textbf{\emph{Pipeline}}: ele
agrega vários processamentos e o estimador final em um único bloco que se comporta
como um estimador (\texttt{fit}/\texttt{predict}), garantindo que cada etapa veja
apenas os dados corretos.

%==============================================================================
\section{O \texttt{StandardScaler} (padronização)}
%==============================================================================
Cada coluna $\bm{X}_i$ de uma base é uma amostra de uma variável aleatória $X_i$ com
$\mu_i=\E[X_i]$ e $\sigma_i^2=\Var(X_i)$ (supostos finitos). A \textbf{padronização}
transforma
\begin{equation}\label{eq:std}
  \widetilde{X}_i = \frac{X_i - \mu_i}{\sigma_i},
\end{equation}
de modo que cada covariável passa a ter média $0$ e variância $1$. Na prática,
$\mu_i$ e $\sigma_i$ são desconhecidos e \emph{estimados} pela média e desvio
amostrais --- e é justamente esse ``estimados'' que cria o risco de vazamento.

\subsection{Vantagens}
\begin{itemize}[leftmargin=1.4cm,nosep]
  \item Covariáveis tornam-se \textbf{adimensionais}: somem as unidades (kg, R\$,
        anos).
  \item Resolve \textbf{discrepâncias numéricas} entre escalas --- crucial para
        métodos baseados em distância (KNN, Aula~04) e em penalização
        (Lasso/Ridge, Aula~02), que de outro modo seriam dominados pelas variáveis
        de maior magnitude.
  \item Torna o \textbf{intercepto desnecessário} em regressão linear (dados
        centrados em zero).
\end{itemize}

\subsection{Desvantagens}
\begin{itemize}[leftmargin=1.4cm,nosep]
  \item Leve perda de \textbf{interpretabilidade}: os coeficientes passam a estar
        em ``desvios-padrão'', não nas unidades originais.
  \item $\mu_i$ e $\sigma_i$ são \textbf{aprendidos dos dados} --- exigindo
        cuidado para não vazar informação do teste (próxima seção).
\end{itemize}

%==============================================================================
\section{Vazamento de dados (\emph{data leakage})}
%==============================================================================
\begin{atencao}[A armadilha]
Se você aplicar o \texttt{StandardScaler} (ou qualquer transformação que aprende
parâmetros) no \emph{conjunto inteiro antes} de separar treino e teste, a
média/desvio usados no treino \textbf{já contêm informação do teste}. As métricas
deixam de medir o desempenho em dados \emph{verdadeiramente novos} e ficam
otimistas. Isso é \textbf{vazamento de dados}.
\end{atencao}

A regra correta:
\begin{enumerate}[leftmargin=1.6cm,nosep]
  \item separe treino e teste \emph{primeiro};
  \item \texttt{fit} do scaler \textbf{só no treino} (aprende $\widehat\mu$,
        $\widehat\sigma$ do treino);
  \item \texttt{transform} aplicado ao treino \emph{e} ao teste com \emph{esses
        mesmos} parâmetros.
\end{enumerate}

\begin{atencao}[Vazamento dentro da validação cruzada]
O mesmo vale para \emph{cada dobra} da CV (Aula~03): o scaler deve ser reajustado
usando apenas o treino \emph{daquela dobra}. Padronizar antes da CV vaza o
``treino futuro'' de cada dobra. Fazer isso manualmente é trabalhoso e propenso a
erro --- por isso usamos \emph{pipelines}, que cuidam disso automaticamente.
\end{atencao}

%==============================================================================
\section{\emph{Pipelines}}
%==============================================================================
Um \texttt{Pipeline} encadeia etapas $(\text{nome}, \text{transformador})$
terminadas por um estimador. Ao chamar:
\begin{itemize}[leftmargin=1.4cm,nosep]
  \item \texttt{pipe.fit(X\_tr, y\_tr)}: cada transformador faz \texttt{fit\_transform}
        \emph{no treino}, em sequência, e o estimador final é ajustado;
  \item \texttt{pipe.predict(X\_te)}: cada transformador apenas \texttt{transform}
        (com os parâmetros aprendidos no treino) e o estimador prediz.
\end{itemize}
Assim, é \emph{impossível} vazar o teste por construção: o teste nunca participa de
nenhum \texttt{fit}.

\begin{lstlisting}
from sklearn.preprocessing import StandardScaler
from sklearn.linear_model import Lasso
from sklearn.pipeline import Pipeline
from sklearn.model_selection import train_test_split, GridSearchCV

X_tr, X_te, y_tr, y_te = train_test_split(X, y, test_size=0.3, random_state=0)

pipe = Pipeline(steps=[("scaler", StandardScaler()),
                       ("lasso",  Lasso())])
pipe.fit(X_tr, y_tr)        # scaler aprende mu/sigma SO do treino
pipe.predict(X_te)          # teste e apenas transformado
\end{lstlisting}

\subsection{\emph{Pipeline} $+$ validação cruzada $+$ busca de hiperparâmetros}
O grande payoff: passar o \emph{pipeline} inteiro ao \texttt{GridSearchCV}. Em cada
dobra, todo o pré-processamento é reajustado corretamente, e a busca pelos
hiperparâmetros (inclusive do scaler, se houver) fica livre de vazamento. Note a
sintaxe \texttt{nomeEtapa\_\_parametro} para acessar parâmetros de cada etapa.

\begin{lstlisting}
from sklearn.linear_model import ElasticNet

pipe_en = Pipeline([("scaler", StandardScaler()),
                    ("enet",   ElasticNet())])

param_grid = {"enet__alpha":    [0.001, 0.01, 0.1, 1, 5, 10, 50, 100],
              "enet__l1_ratio": [0.1, 0.5, 0.7, 0.9, 0.95, 0.99, 1]}

busca = GridSearchCV(pipe_en, param_grid, cv=5,
                     scoring="neg_mean_squared_error")
busca.fit(X_tr, y_tr)               # CV SEM vazamento, etapa por etapa
print("melhores parametros:", busca.best_params_)
print("risco no teste:", -busca.score(X_te, y_te))
\end{lstlisting}

\begin{ideia}
Encare o \emph{pipeline} como ``o modelo''. Tudo o que você faria errado à mão ---
padronizar cedo demais, esquecer de aplicar a mesma transformação no teste, vazar
na CV --- o \emph{pipeline} previne por construção. É a tradução em código das regras
das Aulas~02--04.
\end{ideia}

O notebook \texttt{Aula prática 06} mede os três vazamentos --- seleção de
variáveis, padronização e observações agrupadas --- e monta o
\texttt{ColumnTransformer} completo num banco com faltantes e categóricas.

%==============================================================================
\section{Outros pré-processamentos comuns}
%==============================================================================
\begin{description}[leftmargin=2.6cm,style=nextline]
  \item[Imputação] preencher valores faltantes (\texttt{SimpleImputer} com média,
        mediana ou moda). A estatística de imputação também se aprende \emph{só no
        treino}.
  \item[Codificação de categóricas] \texttt{OneHotEncoder} (cria \emph{dummies}) ou
        \texttt{OrdinalEncoder}. Lembre (Aula~05) que árvores lidam com categóricas
        sem isso.
  \item[Escalas alternativas] \texttt{MinMaxScaler} (intervalo $[0,1]$),
        \texttt{RobustScaler} (baseado em quantis, resistente a \emph{outliers}).
  \item[Redução de dimensão] PCA como etapa do \emph{pipeline} (Aula extra),
        também ajustada apenas no treino.
\end{description}

Quando diferentes colunas precisam de tratamentos diferentes (numéricas
padronizadas, categóricas codificadas), usa-se o \texttt{ColumnTransformer}:
\begin{lstlisting}
from sklearn.compose import ColumnTransformer
from sklearn.preprocessing import OneHotEncoder, StandardScaler
from sklearn.impute import SimpleImputer

prep = ColumnTransformer([
    ("num", Pipeline([("imp", SimpleImputer(strategy="median")),
                      ("sc",  StandardScaler())]), colunas_numericas),
    ("cat", OneHotEncoder(handle_unknown="ignore"), colunas_categoricas),
])
modelo = Pipeline([("prep", prep), ("clf", Lasso())])
\end{lstlisting}

%==============================================================================
\section*{Resumo}
%==============================================================================
\begin{itemize}[leftmargin=1.4cm]
  \item Pré-processamento que aprende parâmetros é \textbf{parte do modelo}: ajuste
        só no treino, aplique ao teste.
  \item \textbf{Padronização} \eqref{eq:std}: adimensionaliza, equaliza escalas
        (essencial para KNN e penalização), dispensa intercepto; custo: leve perda
        de interpretabilidade.
  \item \textbf{Vazamento}: padronizar/imputar antes de separar treino/teste (ou
        antes de cada dobra) infla as métricas. Evite!
  \item \textbf{\emph{Pipeline}}: encadeia tudo num único objeto e previne
        vazamento por construção; combine com \texttt{GridSearchCV}.
  \item \texttt{ColumnTransformer} aplica tratamentos distintos a colunas distintas.
\end{itemize}

\paragraph{Leitura recomendada.} [ISLP] laboratórios dos Capítulos~5 e~6 (uso de
\emph{pipelines} e validação cruzada em Python); [AME] §2.1.2 (regressão linear na
prática). Documentação do \texttt{scikit-learn}: \emph{Pipelines and composite
estimators}.

\end{document}
